\newpage
\Opensolutionfile{ans}[ans/ansCD2D1-2-3]
\begin{dang}{Xác định tham số $m$ để hàm số có cực trị thỏa điều kiện cho trước}
\noindent \textbf{Bài toán 1.} Cho hàm số $y=f(x,m)=ax^3+bx^2+cx+d$. Tìm tham số $m$ để đồ thị hàm số có $2$ điểm cực trị $x_1$, $x_2$ thỏa mãn điều kiện $K$ cho trước?
\begin{enumerate}[\bf Bước 1.]
	\item Tập xác định $\mathscr{D}=\mathbb{R}$. Tính đạo hàm: $y'=3ax^2+2bx+c$.
	\item Hàm số có $2$ điểm cực trị $\Leftrightarrow y'=0$ có $2$ nghiệm phân biệt
	$\Leftrightarrow\heva{&a\neq 0\\&\Delta'=b^2-3ac>0}$ và giải hệ này sẽ tìm được $m\in D_1$.
	\item Gọi $x_1, x_2$ là $2$ nghiệm của phương trình $y'=0$.\\
	Theo định lý Viète, ta có $S=x_1+x_2=-\dfrac{2b}{3a}$ và $P=x_1x_2=\dfrac{c}{3a}$.
	\item Biến đổi điều kiện $K$ về dạng $S$ và $P$. Từ đó giải ra tìm được $m\in D_2$.
	\item Kết luận các giá trị $m$ thỏa mãn: $m=D_1\cap D_2$.
\end{enumerate}
\textbf{Bài toán 2.} Hàm số bậc bốn trùng phương $y=f(x,m)=ax^4+bx^2+c$ có $3$ điểm cực trị thỏa điều kiện $K$.
\begin{enumerate}[\bf Bước 1.]
	\item $y'=4ax^3+2bx=0 \Leftrightarrow \heva{&x=0 \\ & g(x)=2ax^2+b=0.}$
	\item Hàm số có $3$ điểm cực trị $\Leftrightarrow ab<0 \Rightarrow m\in D_1$.
	\item Giải $g(x)=0\Leftrightarrow x_{1,2}=\pm\sqrt{-\dfrac{b}{2a}}\Rightarrow y_1=y_2=f\left(\sqrt{-\dfrac{b}{2a}}\right)$.\\
	Do đó tọa độ ba điểm cực trị sẽ là $A(0;c), B(x_1;y_1), C(x_2;y_2)$ và do tính đối xứng nên tam giác $ABC$ luôn cân tại $A$.
	\item Dựa vào điều kiện đề bài cho để tìm $m\in D_2\Rightarrow m=D_1\cap D_2$.
\end{enumerate}
\end{dang}
\paragraph{Các ví dụ}
\begin{vd}%[Danh Trần]%[2D1K2-5]
Tìm tham số m để các hàm số
\begin{enumerate}
	\item $y=\dfrac{2}{3}x^3-mx^2-2(3m^2-1)x+\dfrac{2}{3}$ có hai điểm cực trị $x_1, x_2$ thỏa mãn $x_1x_2+2(x_1+x_2)=1$.
	\item $y=x^3-3mx+1$ với đồ thị $(C)$ có hai điểm cực trị $B$ và $C$ sao cho tam giác ABC cân tại $A(2;3)$.
	\item $y=2x^3-3(m+1)x^2+6mx$ với đồ thị $(C)$ có hai điểm cực trị $A$ và $B$ sao cho $AB\perp d$, với $d\colon y=x+2$.
	\item $y=x^3-3mx^2+3m^3$ với đồ thị $(C)$ có hai điểm cực trị $A$ và $B$ sao cho $S_{OAB}=48$.
	\item $y=x^4-2(m+1)x^2+m^2$ có $3$ điểm cực trị tạo thành $3$ đỉnh của một tam giác vuông.
	\item $y=x^4-2mx^2$ có ba điểm cực trị tạo thành tam giác có diện tích nhỏ hơn $1$.
\end{enumerate}
\end{vd}
\begin{vd}%[Danh Trần]%[2D1K2-4]
Tìm $m$ để hàm số $y=\dfrac{1}{3}x^3-mx^2+\left(m^2+m-1\right)x+1$ đạt cực trị tại 2 điểm $x_1; x_2$ thỏa mãn $|x_1+x_2|=4$. 
\choice
{$m=2$}
{Không tồn tại $m$}
{\True $m=-2$}
{$m=\pm 2$}
\loigiai{
Ta có $y'=x^2-2mx+m^2+m-1$.\\
Hàm số có hai điểm cực trị khi và chỉ khi $y'=0$ có hai nghiệm phân biệt hay
\[\Delta>0\Leftrightarrow m^2-\left(m^2+m-1\right)>0\Leftrightarrow m<1.\]
Do $x_1$ và $x_2$ là nghiệm của phương trình $y'=0$ nên $x_1+x_2=2m$.\\
Theo giả thiết $|x_1+x_2|=4\Leftrightarrow|2m|=4\Leftrightarrow\hoac{&m=2\\&m=-2}$. So với điều kiện, ta nhận $m=-2$.}
\end{vd}

\paragraph{Câu hỏi trắc nghiệm}
\begin{ex}%[Danh Trần]%[2D1K2-4]
	Cho hàm số $y=-x^3+(2m+1)x^2-\left(m^2-3m+2\right)x-4$. Tìm $m$ để hàm số có cực đại, cực tiểu nằm 2 phía trục tung
	\choice
	{\True $m\in(1;2)$}
	{$m\in[1;2]$}
	{$m\in(-\infty;1)\cup(2;+\infty)$}
	{$m\in(-\infty;1]\cup[2;+\infty)$}
	\loigiai{
		Đạo hàm $y'=-3x^2+2(2m+1)x-(m^2-3m+2)$. Gọi $x_1$ và $x_2$ là 2 cực trị của hàm số.\\
		Theo yêu cầu bài toán ta có $x_1\cdot x_2<0\Leftrightarrow\dfrac{m^2-3m+2}{3}<0\Leftrightarrow 1<m<2$.
	}
\end{ex}
\begin{ex}%[Danh Trần]%[2D1K2-4]
	Tìm tất cả các giá trị của tham số $m$ để hàm số $y=\dfrac{1}{3}x^3-\dfrac{1}{2}(m+5)x^2+mx$ có điểm cực đại, cực tiểu và $\left|x_{\text{CĐ}}-x_{\text{CT}}\right|=5$. 
	\choice
	{$m=0$}
	{\True $m=\{-6;0\}$}
	{$m=6$}
	{$m=\{6;0\}$}
	\loigiai{
		Ta có $y'=x^2-(m+5)x+m$.\\
		Để hàm số có điểm cực đại và cực tiểu thì $y'=x^2-(m+5)x+m=0$ có hai nghiệm phân biệt.\\
		Suy ra $\Delta>0\Leftrightarrow(m+5)^2-4m>0\Leftrightarrow m^2+6m+25>0$, $\forall m\in\mathbb{R}$.\\
		Áp dụng định lý Vi-ét ta có
		$\heva{&x_{\text{CĐ}}+x_{\text{CT}}=m+5\\&x_{text{CĐ}}\cdot x_{\text{CT}}=m}$. Khi đó \[\left|x_{\text{CĐ}}-x_{\text{CT}}\right|=5\Leftrightarrow\left(x_{\text{CĐ}}+x_{\text{CT}}\right)^2-4x_{C\S}\cdot x_{\text{CT}}=25\Leftrightarrow(m+5)^2-4m-25=0\Leftrightarrow\hoac{&m=0\\&m=-6.}\]
	}
\end{ex}
\begin{ex}%[Danh Trần]%[2D1K2-4]
	Gọi $x_1,x_2$ là hai điểm cực trị của hàm số $y=x^3-3mx^2+3\left(m^2-1\right)x-m^3+m$. Tìm tất cả các giá trị của tham số thực $m$ để $x_1^2+x_2^2-x_1x_2=7$. 
	\choice
	{$m=\pm 1$}
	{\True $m=\pm 2$}
	{$m=0$}
	{$m=\pm\sqrt{2}$}
	\loigiai{
		$y'=3x^2-6mx+3\left(m^2-1\right)$; $\Delta'=9>0,\forall m$ nên hàm số luôn có 2 cực trị $x_1,x_2$ thỏa $\heva{&x_1+x_2=2m\\&x_1x_2=m^2-1.}$\\Khi đó $x_1^2+x_2^2-x_1x_2=7\Leftrightarrow(x_1+x_2)^2-3x_1x_2=7\Leftrightarrow m^2=4\Leftrightarrow m=\pm 2$.}
\end{ex}
\begin{ex}%[Danh Trần]%[2D1K2-4]
	Đồ thị hàm số $y=x^3-(3m+1)x^2+\left(m^2+3m+2\right)x+3$ có điểm cực tiểu và điểm cực đại nằm về hai phía của trục tung khi: 
	\choice
	{$1<m<2$}
	{\True $-2<m <-1$}
	{$2<m<3$}
	{$-3<m <-2$}
	\loigiai{
		Ta có $y'=3x^2-2(3m+1)x+m^2+3m+2$. Yêu cầu bài toán tương đương với \[ac<0\Leftrightarrow 3\cdot\left(m^2+3m+2\right)<0\Leftrightarrow\left(m^2+3m+2\right)<0\Leftrightarrow-2<m <-1.\]}
\end{ex}
\begin{ex}%[Danh Trần]%[2D1K2-4]
	Với giá trị nào của $m$ thì đồ thị hàm số $y=\dfrac{1}{3}x^3-mx^2+(m+2)x$ có hai điểm cực trị nằm về phía bên phải trục tung?
	\choice
	{$m=2$}
	{$0<m<2$}
	{$m<2$}
	{\True $m>2$}
	\loigiai{
		Ta có tập xác định $\mathscr{D}=\mathbb{R}$, đạo hàm $y'=x^2-2mx+m+2$.\\
		Đồ thị hàm số có hai điểm cực trị nằm về bên phải trục tung khi phương trình $y'=0$ có hai nghiệm dương phân biệt, khi đó
		$\heva{&m^2-m-2>0\\&2m>0\\&m+2>0}\Leftrightarrow\heva{&\hoac{&m <-1\\&m>2}\\&m>0\\&m >-2}\Leftrightarrow m>2$.
	}
\end{ex}
\begin{ex}%[Danh Trần]%[2D1K2-4]
	Tìm $m$ để hàm số $y=\dfrac{1}{3}x^3-mx^2+(m+2)x-1$ có hai điểm cực trị trong khoảng $(0;+\infty)$?
	\choice
	{$(-\infty;2]$}
	{\True $(2;+\infty)$}
	{$(-2;4)$}
	{$[2;+\infty)$}
	\loigiai{
		Ta có: $y'=x^2-2mx+m+2$.\\
		Hàm số có 2 điểm cực trị trong $(0;+\infty)$ khi và chỉ khi phương trình $y'=0$ có 2 nghiệm dương phân biệt hay $\Leftrightarrow\heva{&\Delta'>0\\&P>0\\&S>0}\Leftrightarrow\heva{&m^2-m-2>0\\&2m>0\\&m+2>0}\Leftrightarrow\heva{&m <-1\vee m>2\\&m>0\\&m >-2}\Leftrightarrow m>2$.}
\end{ex}
\begin{ex}%[Danh Trần]%[2D1K2-4]
	Tìm tất cả các giá trị thực của tham số $a$ sao cho hàm số $y=\dfrac{1}{3}x^3-\dfrac{1}{2}x^2+ax+1$ đạt cực trị tại điểm $x_1,x_2$ thỏa mãn: $\left(x_1^2+x_2+2a\right)\left(x_2^2+x_1+2a\right)=9$ 
	\choice
	{$a=2$}
	{\True $a=-4$}
	{$a=-3$}
	{$a=-1$}
	\loigiai{
		Đạo hàm $y'=x^2-x+a$. Xét phương trình $y'=0\Leftrightarrow x^2-x+a=0$.\\
		Để hàm số có hai cực trị $\Leftrightarrow y'=0$ có 2 nghiệm phân biệt $\Leftrightarrow 1-4a>0\Leftrightarrow a<\dfrac{1}{4}$.\\
		Ta có theo định lý Vi-ét $S=x_1+x_2=1$ và $P=x_1x_2=a$.\\
		Ta có $x_1^2-x_1+a=0\Leftrightarrow x_1^2=x_1-a$ và $x_2^2-x_2+a=0\Leftrightarrow x_2^2=x_2-a$.\\
		Theo đề bài, $\left(x_1^2+x_2+2a\right)\left(x_2^2+x_1+2a\right)=9$
		\[\Leftrightarrow\left(x_1-a+x_2+2a\right)\left(x_2-a+x_1+2a\right)=9\Leftrightarrow(x_1+x_2+a)^2=9\Leftrightarrow(1+a)^2=9\Leftrightarrow\hoac{&a=2\\&a=-4.}\]
		So sánh điều kiện, nhận $a=-4$.
	}
\end{ex}
\begin{ex}%[Danh Trần]%[2D1K2-4]
	Cho hàm số $y=\dfrac{2}{3}x^3+(m+1)x^2+\left(m^2+4m+3\right)x-3$, ($m$ là tham số thực). Tìm điều kiện của $m$ để hàm số có cực đại cực tiểu và các điểm cực trị của đồ thị hàm số nằm bên phải của trục tung. 
	\choice
	{$-5<m <-1$}
	{\True $-5<m <-3$}
	{$-3<m <-1$}
	{$\hoac{&m >-1\\&m <-5}$}
	\loigiai{
		$y'=2x^2+2(m+1)x+\left(m^2+4m+3\right)$.\\
		Yêu cầu bài toán thỏa mãn trở thành $y'=0$ có hai nghiệm dương phân biệt, từ đó \[\heva{&\Delta'>0\\&S>0\\&P>0}\Leftrightarrow\heva{&(m+1)^2-2\left(m^2+4m+3\right)>0\\&-(m+1)>0\\&\dfrac{m^2+4m+3}{2}>0}\Leftrightarrow\heva{&m\in(-5;-1)\\&m <-1\\&m\in(-\infty;-3)\cup(-1;+\infty)}\Leftrightarrow m\in(-5;-3)\]
		Vậy $-5<m<3$ thỏa mãn yêu cầu bài toán.
	}
\end{ex}
\begin{ex}%[Danh Trần]%[2D1K2-4]
	Cho hàm số $y=x^3-3x^2+mx-1$, tìm giá trị của tham số $m$ để hàm số có hai cực trị $x_1$, $x_2$ thỏa $x_1^2+x_2^2=3$. 
	\choice
	{\True $m=\dfrac{3}{2}$}
	{$m=1$}
	{$m=-2$}
	{$m=\dfrac{1}{2}$}
	\loigiai{
		Tập xác định $\mathscr{D}=\mathbb{R}$. Ta có đạo hàm $f'(x)=3x^2-6x+m$.\\
		Hàm số có hai cực trị $x_1,x_2$ khi $f'(x)=0$ có hai nghiệm phân biệt $\Leftrightarrow 9-3m>0\Leftrightarrow m<3$.\\
		Theo hệ thức Vi-et, $x_1+x_2=2$ và $x_1\cdot x_2=\dfrac{m}{3}$.\\
		Ta có: $x_1^2+x_2^2=3\Leftrightarrow(x_1+x_2)^2-2x_1x_2=3\Leftrightarrow 2^2-2\cdot\dfrac{m}{3}=3\Leftrightarrow m=\dfrac{3}{2}$.}
\end{ex}
\begin{ex}%[Danh Trần]%[2D1K2-5]
	Cho hàm số $y=3x^4-2mx^2+2m+m^4$. Tìm tất cả các giá trị của $m$ để đồ thị hàm số đã cho có ba điểm cực trị tạo thành tam giác có diện tích bằng $3$. 
	\choice
	{$m=-3$}
	{\True $m=3$}
	{$m=4$}
	{$m=-4$}
	\loigiai{
		Ta có $y'=12x^3-4mx =4x\left(3x^2-m\right)$. Đề đồ thị hàm số có ba điểm cực trị thì $m>0$.\\
		Khi đó tọa độ các điểm cực trị là $A\left(0;2m+m^4\right)$, $B\left(\sqrt{\dfrac{m}{3}};m^4-\dfrac{m^2}{3}+2m\right)$, $C\left(-\sqrt{\dfrac{m}{3}};m^4-\dfrac{m^2}{3}+2m\right)$.\\
		Tam giác $ABC$ cân tại $A$ nên có diện tích $S_{ABC}=\dfrac{1}{2}\cdot BC\cdot\mathrm{d}(A;BC) =\dfrac{1}{2}\cdot 2\sqrt{\dfrac{m}{3}}\cdot\dfrac{m^2}{3} =\sqrt{\dfrac{m}{3}}\cdot\dfrac{m^2}{3}$.\\
		Theo đề bài ta có $\sqrt{\dfrac{m}{3}}\cdot\dfrac{m^2}{3}=3\Leftrightarrow m=3$.}
\end{ex}
\begin{ex}%[Danh Trần]%[2D1K2-4]
	Tìm $m$ để hàm số $y=x^3-3x^2+mx-m^3$ có hai điểm cực trị $x_1$; $x_2$ thỏa mãn $x_1^2+x_2^2=3$. 
	\choice
	{$m=-\dfrac{3}{2}$}
	{$m=-3$}
	{$m=3$}
	{\True $m=\dfrac{3}{2}$}
	\loigiai{
		Tập xác định $\mathscr{D}=\mathbb{R}$. Ta có $y'=3x^2-6x+m$.\\
		Vậy hàm số có hai điểm cực trị $x_1$; $x_2$ thỏa mãn $x_1^2+x_2^2=3$ khi và chỉ khi $y'=0$ có hai nghiệm phân biệt $x_1$; $x_2$ thỏa mãn $x_1^2+x_2^2=3\Leftrightarrow\heva{&\Delta'>0\\&x_1^2+x_2^2=3}\Leftrightarrow\heva{&9-3m>0\\&(x_1+x_2)^2-2x_1x_2=3}\Leftrightarrow\heva{&m<3\\&2m=3}\Leftrightarrow m=\dfrac{3}{2}$.}
\end{ex}
\begin{ex}%[Danh Trần]%[2D1K2-4]
	Cho hàm số $y=\dfrac{1}{3}x^3-(m+2)x^2+\left(m^2+4m+3\right)+5m^3+1(1)$. Gọi $m$ là số thực để hàm số (1) đạt cực đại tại $x_1$, đạt cực tiểu tại $x_2$ sao cho $x_1^2=x_2$. Khi đó khẳng định nào sau đây đúng. 
	\choice
	{\True $m\in(-3;3)$}
	{$m>6$}
	{$m\in(3;6)$}
	{$m <-1$}
	\loigiai{
		Ta có $y'=x^2-2(m+2)x+m^2+4m+3$. Tại $y'=0\Leftrightarrow x^2-2(m+2)x+m^2+4m+3=0\Leftrightarrow\hoac{&x=m+1\\&x=m+3.}$ \\
		Ta có biệt thức $\Delta'=(m+2)^2-\left(m^2+4m+3\right)=1>0$. Do đó hàm số luôn có 2 điểm cực trị.
		\immini
		{
		Bảng biến thiên như hình bên.\\
		Khi đó $x_1=m+1$, $x_2=m+3$.\\
		Do $x_1^2=x_2$ nên \[(m+1)^2=m+3\Leftrightarrow m^2+m-2=0\Leftrightarrow\hoac{&m=1\\&m=-2.}\]
		}
		{
		\begin{tikzpicture}[scale=0.9]
		\tkzTabInit[nocadre=false,lgt=1.2,espcl=2.5,deltacl=0.6]
		{$x$/1, $y'$/1, $y$/2}
		{$-\infty$, $m+1$, $m+3$, $+\infty$}
		\tkzTabLine{,+,0,-,0,+,}
		\tkzTabVar{-/$-\infty$, +/ $y(m+1)$, -/$y(m+3)$, +/ $+\infty$}
		\end{tikzpicture}
		}
	}
\end{ex}
\begin{ex}%[Danh Trần]%[2D1K2-5]
	Tìm tất cả giá trị của tham số $m$ để đồ thị hàm số $y=-x^4+2mx^2$ có ba điểm cực trị tạo thành một tam giác đều. 
	\choice
	{$m>0$}
	{\True $m=\sqrt[3]{3}$}
	{$m=\pm\sqrt[3]{3}$}
	{$m=1$}
	\loigiai{
		Tập xác định $\mathscr{D}=\mathbb{R}$. Đạo hàm $y'=-4x^3+4mx$. Tại $y'=0\Leftrightarrow\hoac{&x=0\\&x^2=m.}$\\
		Hàm số có 3 điểm cực trị $\Leftrightarrow y'=0$ có 3 nghiệm phân biệt $\Leftrightarrow m>0$.\\
		Lúc đó đồ thị hàm số có 3 điểm cực trị $A(0;0),B\left(-\sqrt{m};m^2\right),C\left(\sqrt{m};m^2\right)$.\\
		Ta có $AB=AC=\sqrt{m^4+m},BC=2\sqrt{m}$.\\
		Theo giả thiết, ta có $AB=BC\Leftrightarrow\sqrt{m+m^4}=2\sqrt{m}\Leftrightarrow m^4-3m=0\Leftrightarrow\hoac{&m=0\\&m=\sqrt[3]{3}}\Leftrightarrow m=\sqrt[3]{3}$.}
\end{ex}
\begin{ex}%[Danh Trần]%[2D1K2-5]
	Gọi $m_0$ là giá trị thực của tham số $m$ để đồ thị hàm số $y=x^4+2mx^2+4$ có 3 điểm cực trị nằm trên các trục tọa độ. Mệnh đề nào sau đây là đúng?
	\choice
	{$m_0\in(-5;-3)$}
	{$m_0\in\left(-\dfrac{3}{2};0\right)$}
	{\True $m_0\in\left(-3;-\dfrac{3}{2}\right)$}
	{$m_0\in(1;3)$}
	\loigiai{
		Ta có $y=x^4+2mx^2+4\Rightarrow y'=4x^3+4mx$.\\
		Với $y'=0\Leftrightarrow 4x^3+4mx=0\Leftrightarrow\hoac{&x=0\Rightarrow y=4\\&x=\sqrt{-m}\Rightarrow y=4-m^2\\&x=-\sqrt{-m}\Rightarrow y=4-m^2}$ ($m<0$).\\
		Yêu cầu bài toán xảy ra khi $4-m^2=0\Leftrightarrow m=\pm 2\Rightarrow m=-2 (m<0)$.}
\end{ex}
\begin{ex}%[Danh Trần]%[2D1K2-4]
	Cho hàm số $y=-\dfrac{1}{3}x^3+x^2+mx+1$. Tìm tập hợp các giá trị của $m$ để hàm số đạt cực trị tại các điểm $x_1,x_2$ thỏa mãn $x_1^2+x_2^2=6$. 
	\choice
	{$\{0\}$}
	{$(-1;+\infty)$}
	{$\{2\}$}
	{\True $\{1\}$}
	\loigiai{
		Đạo hàm $y'=-x^2+2x+m$.\\
		Để hàm số có hai điểm cực trị thì $y'=0$ có hai nghiệm phân biệt $\Leftrightarrow\Delta'=1+m>0\Leftrightarrow m >-1$.\\
		Khi $m >-1$, ta có hoành độ cực trị $x_1,x_2$ là nghiệm của phương trình $y'=0$.\\
		Theo định lý Vi-ét ta có $\heva{&x_1+x_2=2\\&x_1\cdot x_2=-m}$. Yêu cầu bài toán là
		\[x_1^2+x_2^2=6\Leftrightarrow(x_1+x_2)^2-2x_1\cdot x_2=6\Leftrightarrow 4+2m=6\Leftrightarrow m=1.\]}
\end{ex}
\begin{ex}%[Danh Trần]%[2D1K2-4]
	Cho đồ thị hàm số $y=-x^3+3mx+1$ có hai điểm cực trị $A,B$ thỏa mãn tam giác $OAB$ vuông tại $O$ ($O$ là gốc tọa độ). Khẳng định nào dưới đây đúng?
	\choice
	{$-1<m<\dfrac{1}{3}$}
	{$1<m<3$}
	{\True $-\dfrac{1}{2}<m<1$}
	{$-2<m<0$}
	\loigiai{
		Đạo hàm$y'=-3x^2+3m$. Tại $y'=0\Leftrightarrow-3x^2+3m=0\Leftrightarrow x^2=m\Leftrightarrow m>0$.\\
		Khi $m>0$, phương trình $y'=0$ có hai nghiệm phân biệt là $x=\pm\sqrt{m}$.\\
		Gọi $A\left(\sqrt{m};1+2m\sqrt{m}\right),B\left(\sqrt{m};1-2m\sqrt{m}\right)$ là hai điểm cực trị của $(C_m)$.\\
		Để tam giác $OAB$ vuông tại $O$ thì
		\[\overrightarrow{OA}\cdot\overrightarrow{OB}=0\Leftrightarrow\sqrt{m}\cdot(-\sqrt{m})+\left(1+2m\sqrt{m}\right)\left(1-2m\sqrt{m}\right)=0\Leftrightarrow-4m^3-m+1=0\Leftrightarrow m=\dfrac{1}{2}.\]}
\end{ex}
\begin{ex}%[Danh Trần]%[2D1K2-5]
	Tìm tất cả các giá trị thực của tham số $m$ để đồ thị của hàm số $y=x^4-2mx^2$ có ba điểm cực trị tạo thành một tam giác có diện tích nhỏ hơn $1$. 
	\choice
	{$m<1$}
	{\True $0<m<1$}
	{$0<m<\sqrt[3]{4}$}
	{$m>0$}
	\loigiai{
		Hàm số $y=x^4-2mx^2$ có. Ta có $y'=4x^3-4mx$; $y'=0\Leftrightarrow\hoac{&x=0\\&x^2=m.}$ \\
		Để đồ thị hàm số có ba điểm cực trị thì $m>0$.\\
		Khi đó ba điểm cực trị là $O(0;0)$, $B\left(-\sqrt{m};-m^2\right)$, $C\left(\sqrt{m};-m^2\right)$.\\
		Ta có tam giác $OBC$ cân tại $O$, với $I\left(0;-m^2\right)$ là trung điểm của $BC$.\\
		Theo yêu cầu bài toán, ta có $S_{ABC}=\dfrac{1}{2}OI\cdot BC=\dfrac{1}{2}|-m^2|\cdot 2\sqrt{m}<1\Leftrightarrow 0<m<1$.}
\end{ex}
\begin{ex}%[Danh Trần]%[2D1K2-4]
	Điểm $M(3;-1)$ thuộc đường thẳng đi qua hai điểm cực đại và cực tiểu của đồ thị hàm số $y=x^3-x+m$ khi $m$ bằng
	\choice
	{$2$}
	{\True $1$}
	{$-1$}
	{$0$}
	\loigiai{
		Ta có $y'=3x^2-1$; $\dfrac{y}{y'}=\dfrac{x^3-x+m}{3x^2-1}=\dfrac{1}{3}x+\dfrac{-\dfrac{2}{3}x+m}{3x^2-1}=\dfrac{1}{3}x\left(3x^2-1\right)-\dfrac{2}{3}x+m$.\\
		Suy ra phương trình qua 2 điểm cực đại và cực tiểu là $y=\dfrac{-2}{3}x+m$.\\
		Thay $M(3;-1)$ vào ta có $-1=\dfrac{-2}{3}\cdot3+m\Rightarrow-1=-2+m\Rightarrow m=1$.}
\end{ex}
\begin{ex}%[Danh Trần]%[2D1K2-4]
	Giả sử đồ thị của hàm số $y=x^3-3mx^2+m$ có hai điểm cực trị $A,B$. Tìm tất cả các giá trị thực của tham số $m$ để đường thẳng $AB$ song song với đường thẳng $d\colon y=1-2x$?
	\choice
	{$m=1$}
	{\True $m=-1$}
	{$m=1$; $m=-1$}
	{$m=2$; $m=-2$}
	\loigiai{
		Công thức PTĐT qua điểm cực trị là $y=\dfrac{-2\left(b^2-3ac\right)}{9a}x+\dfrac{9ad-bc}{9a}$.\\
		PTĐT qua hai điểm cực trị $A,B$ là $AB\colon y=-2m^2x+m$.\\
		Vì $AB\parallel d$ nên $\heva{&-2m^2=-2\\&m\neq 1}\Leftrightarrow\heva{&\hoac{&m=1\\&m=-1}\\&m\neq 1}\Leftrightarrow m=-1$.}
\end{ex}
\begin{ex}%[Danh Trần]%[2D1K2-4]
	Tìm tất cả các giá trị thực của tham số $m$ để đồ thị hàm số $y=x^3-3mx^2+4m^3$ có hai điểm cực trị $A$ và $B$ sao cho tam giác $OAB$ có diện tích bằng $4$ với $O$ là gốc tọa độ. 
	\choice
	{\True $m=-\dfrac{1}{\sqrt{2}}$; $m=\dfrac{1}{\sqrt{2}}$}
	{$m=-1$; $m=1$}
	{$m=1$}
	{$m\neq 0$}
	\loigiai{
		Ta có $y'=3x^2-6mx=3x(x-2m)$ và $y''=6x-6m$. Tại $y'=0\Leftrightarrow\hoac{&x=0\\&x=2m.}$ \\
		Để hàm số có hai điểm cực trị thì $y''(0)\neq 0$ và $y''(2m)\neq 0$, hay $m\neq 0$.\\
		Khi đó, ta có tọa độ hai điểm $A$ và $B$ là $A\left(0; 4m^3\right)$, $B(2m; 0)$.\\
		Vậy tam giác $OAB$ vuông tại $O$ và $S_{OAB}=\dfrac{1}{2}OA\cdot OB=\dfrac{1}{2}\left|4m^3\right|\cdot|2m|=4m^4$.\\
		Ta cần $S_{OAB}=1\Leftrightarrow 4m^4=1\Leftrightarrow m=\dfrac{1}{\sqrt{2}}$ hoặc $m=-\dfrac{1}{\sqrt{2}}$ (thỏa mãn $m\neq 0$).}
\end{ex}
\Closesolutionfile{ans}
